% Options for packages loaded elsewhere
\PassOptionsToPackage{unicode}{hyperref}
\PassOptionsToPackage{hyphens}{url}
\PassOptionsToPackage{dvipsnames,svgnames,x11names}{xcolor}
\documentclass[
  a4paper,
]{article}
\usepackage{xcolor}
\usepackage[margin=1in]{geometry}
\usepackage{amsmath,amssymb}
\setcounter{secnumdepth}{5}
\usepackage{iftex}
\ifPDFTeX
  \usepackage[T1]{fontenc}
  \usepackage[utf8]{inputenc}
  \usepackage{textcomp} % provide euro and other symbols
\else % if luatex or xetex
  \usepackage{unicode-math} % this also loads fontspec
  \defaultfontfeatures{Scale=MatchLowercase}
  \defaultfontfeatures[\rmfamily]{Ligatures=TeX,Scale=1}
\fi
\usepackage{lmodern}
\ifPDFTeX\else
  % xetex/luatex font selection
\fi
% Use upquote if available, for straight quotes in verbatim environments
\IfFileExists{upquote.sty}{\usepackage{upquote}}{}
\IfFileExists{microtype.sty}{% use microtype if available
  \usepackage[]{microtype}
  \UseMicrotypeSet[protrusion]{basicmath} % disable protrusion for tt fonts
}{}
\usepackage{longtable,booktabs,array}
\newcounter{none} % for unnumbered tables
\usepackage{calc} % for calculating minipage widths
% Correct order of tables after \paragraph or \subparagraph
\usepackage{etoolbox}
\makeatletter
\patchcmd\longtable{\par}{\if@noskipsec\mbox{}\fi\par}{}{}
\makeatother
% Allow footnotes in longtable head/foot
\IfFileExists{footnotehyper.sty}{\usepackage{footnotehyper}}{\usepackage{footnote}}
\makesavenoteenv{longtable}
\usepackage{graphicx}
\makeatletter
\newsavebox\pandoc@box
\newcommand*\pandocbounded[1]{% scales image to fit in text height/width
  \sbox\pandoc@box{#1}%
  \Gscale@div\@tempa{\textheight}{\dimexpr\ht\pandoc@box+\dp\pandoc@box\relax}%
  \Gscale@div\@tempb{\linewidth}{\wd\pandoc@box}%
  \ifdim\@tempb\p@<\@tempa\p@\let\@tempa\@tempb\fi% select the smaller of both
  \ifdim\@tempa\p@<\p@\scalebox{\@tempa}{\usebox\pandoc@box}%
  \else\usebox{\pandoc@box}%
  \fi%
}
% Set default figure placement to htbp
\def\fps@figure{htbp}
\makeatother
% definitions for citeproc citations
\NewDocumentCommand\citeproctext{}{}
\NewDocumentCommand\citeproc{mm}{%
  \begingroup\def\citeproctext{#2}\cite{#1}\endgroup}
\makeatletter
 % allow citations to break across lines
 \let\@cite@ofmt\@firstofone
 % avoid brackets around text for \cite:
 \def\@biblabel#1{}
 \def\@cite#1#2{{#1\if@tempswa , #2\fi}}
\makeatother
\newlength{\cslhangindent}
\setlength{\cslhangindent}{1.5em}
\newlength{\csllabelwidth}
\setlength{\csllabelwidth}{3em}
\newenvironment{CSLReferences}[2] % #1 hanging-indent, #2 entry-spacing
 {\begin{list}{}{%
  \setlength{\itemindent}{0pt}
  \setlength{\leftmargin}{0pt}
  \setlength{\parsep}{0pt}
  % turn on hanging indent if param 1 is 1
  \ifodd #1
   \setlength{\leftmargin}{\cslhangindent}
   \setlength{\itemindent}{-1\cslhangindent}
  \fi
  % set entry spacing
  \setlength{\itemsep}{#2\baselineskip}}}
 {\end{list}}
\usepackage{calc}
\newcommand{\CSLBlock}[1]{\hfill\break\parbox[t]{\linewidth}{\strut\ignorespaces#1\strut}}
\newcommand{\CSLLeftMargin}[1]{\parbox[t]{\csllabelwidth}{\strut#1\strut}}
\newcommand{\CSLRightInline}[1]{\parbox[t]{\linewidth - \csllabelwidth}{\strut#1\strut}}
\newcommand{\CSLIndent}[1]{\hspace{\cslhangindent}#1}
\setlength{\emergencystretch}{3em} % prevent overfull lines
\providecommand{\tightlist}{%
  \setlength{\itemsep}{0pt}\setlength{\parskip}{0pt}}
\usepackage{amsmath}
\numberwithin{equation}{section}
\numberwithin{figure}{section}
\numberwithin{table}{section}
\usepackage{setspace}\onehalfspacing
\usepackage{pdflscape}
\newcommand{\blandscape}{\begin{landscape}}
\newcommand{\elandscape}{\end{landscape}}
\usepackage{indentfirst}
\usepackage{xparse}
\usepackage{tcolorbox}
\usepackage{algorithm}
\usepackage{algpseudocode}
\usepackage{float}
\floatplacement{figure}{H}
\usepackage{placeins}
\renewcommand{\ttfamily}{\rmfamily}
\usepackage{etoolbox}
\usepackage{multirow}
\usepackage{makecell}
\usepackage{booktabs}
\AtBeginEnvironment{algorithmic}{\setstretch{1.5}}
\renewcommand{\arraystretch}{0.75}
\renewcommand{\multirowsetup}{\raggedright}
\usepackage{anyfontsize}
\newcommand{\thesisfont}{\fontsize{12pt}{12pt}\selectfont}
\thesisfont
\usepackage{xcolor}
\usepackage{titlesec}
\usepackage{caption}
\captionsetup[figure]{font={large,bf}}
\captionsetup[table]{font={large,bf}}
\newcommand{\algfont}{\fontsize{15pt}{18pt}\selectfont}
\usepackage{bookmark}
\IfFileExists{xurl.sty}{\usepackage{xurl}}{} % add URL line breaks if available
\urlstyle{same}
\hypersetup{
  pdftitle={Modeling Workers' Return-to-Office Decisions in Australia},
  pdfauthor={Shaine Rosewel Matala},
  colorlinks=true,
  linkcolor={blue},
  filecolor={Maroon},
  citecolor={blue},
  urlcolor={blue},
  pdfcreator={LaTeX via pandoc}}

\title{Modeling Workers' Return-to-Office Decisions in Australia}
\author{Shaine Rosewel Matala}
\date{2026-04-23}

\begin{document}
\maketitle
\begin{abstract}
This paper uses the Global Survey of Working Arrangements (G-SWA, 2022) to describe respondents from the working class of Australia based on the decision they would take if their employer announced that all employees must return to the worksite 5+ days a week by a certain date. The response signals a preference for either a hybrid/WFH or 100\% in-office work setting. We aim to describe this decision based on their immediate environment's perception of remote work, commute time, gender, age, whether they have children, industry, and other factors. We consider a generalized linear model with an identity or logistic link. This may be beneficial for employers and the government in understanding workers' needs and potentially informing regulations.
\end{abstract}

\section{Introduction}\label{sec:intro}

Working from home is a formal arrangement where employees perform duties without daily physical attendance at the office, instead doing so remotely from a place of work such as the home or a café, at conventional hours or flexibly at other times, enabled by information and communication technologies (Akshay Vij (\citeproc{ref-vij-2023}{2023})). The COVID-19 pandemic in 2020 paved the way for a significant Work from Home (WfH) uptake across Australia. The Department of Infrastructure, Transport, Regional Development, Communications and the Arts (DITRDCSA) (\citeproc{ref-dit-2022}{2022}) defines WfH uptake as the ``proportion of job tasks and activities that were done remotely at a given point in time'' (p.~7).

In late 2020 and early 2021, few employers believed that the hybrid model would be sustained beyond 2021, and some strongly rejected WfH as a long-term strategy. Meanwhile, others recognized the need to adapt in terms of flexibility, while some viewed hybrid arrangements as appropriate only for certain roles. Managerial employees, on the other hand, were generally more supportive of company policies allowing WfH in the long term (DITRDCSA, 2022). These differing perspectives may translate into varying workplace policies, including mandates to return to the office. However, there is limited understanding of how workers immediately respond to strict return-to-office mandates, particularly in terms of their decision to comply or seek alternative work arrangements.

In this context, this research examines the immediate decision of workers when faced with a mandate to return to the worksite 5+ days a week. In particular, it considers whether workers prefer to remain in a WfH arrangement (e.g., seeking a job that allows remote work or leaving their current job) or to return fully to in-office work (i.e., complying with the mandate). It further aims to understand the factors associated with this decision, including perception of remote work, commute time, gender, age, presence of children, and industry.

Results may be beneficial for employers and policymakers in understanding how workers respond to strict return-to-office mandates. The findings may help inform workplace policies and labor regulations, particularly in balancing organizational requirements with worker preferences and constraints. We use the Global Survey of Working Arrangements (G-SWA, 2022) dataset and focuses on working-class respondents in Australia. It is limited to individuals presented with a hypothetical return-to-office mandate and does not account for actual enforced policy outcomes.

\section{Related Literature}\label{related-literature}

\subsection{Factors Associated to Work Arrangement Preference}\label{factors-associated-to-work-arrangement-preference}

Rüger, Laß, Stawarz, \& Mergener (\citeproc{ref-ruger-2024}{2024}) show in their study employees may be able to save an average of 14\% of their time from commuting. They add that this would increase in the long run. Furthermore, they explain that only high shares of WfH are linked to substantial drops in commuting time, and reductions are larger for women than men. This may be explained by differences in job selection and commuting tolerance, where men with low WFH shares are more likely to accept jobs involving longer commutes if they require fewer days in the office, while women may place relatively higher value on jobs with shorter commuting distances. Overall, when accounting for workers' reported WfH preferences, their results suggest maximum potential future commuting time savings of about 17--25\% compared to 2019.

The 2024 Household Income and Labour Dynamics in Australia (HILDA) Survey data shows workforce participation in jobs allowing WfH increased by 9 percentage points for women with children under 4 years old and 4.4 percentage points for people with a disability from 2019 to 2023. It is also observed that prior to the pandemic, carers, those with limiting condition, and women with children are among those with the highest WfH rates. With the onset of the pandemic until 2023, rate differences between the former group and general workers, as well as women without children decreased, making WfH rates more similar across groups (Committee for Economic Development of Australia (CEDA) (\citeproc{ref-ceda-2024}{2024})). This suggests that even before the pandemic, the former group had a higher tendency toward WfH. As Barrero, Bloom, \& Davis (\citeproc{ref-barrero-2022}{2022}) states, married persons---both men and women who are also carers---tend to place more value on the WfH option when they have children aged 14 years or younger.

In their study, Akshay Vij (\citeproc{ref-vij-2023}{2023}) shows that an average worker would be willing to forego around 4 to 8 percent of their annual wages in exchange of an ability to WfH on some workdays and/or workhours. In terms age, workers in their 20's have the lowest WfH valuation and this is opposite to those in their 30's and 50's. They explain that this is associated to younger workers being at the start of their career and therefore tending to seek direct interaction with colleagues and guidance from supervisors. They also find that more educated workers are more willing to accept a pay cut for an ability to WfH as their jobs require a greater capability of being done remotely. Women also tend to value the flexibility of WfH more than men do, a pattern often attributed in the literature to their greater caregiving and household responsibilities (e.g.~Wiswall and Zafar, 2017). In terms of household structure, people who live alone have the smallest compensating wage differential, indicating that only a small wage premium is required for them to accept working in the office. In contrast, couples with children living at home are willing to forgo the largest share of their salary in exchange for the option to work from home.

The DITRDCSA defines \emph{``WfH uptake as the the proportion of job tasks and activities that were done remotely at a given point in time''}. It is further noted that \emph{``individuals are considered to have WfH capability if they could have done their job remotely on some days in a week if appropriate company policies and resources for remote working were put in place''}. Undeniably, WfH capability differs based on the occupation of the individual. As to the report, WfH capability and uptake are much higher for managerial employees than non-managerial employees. Beck \& Hensher (\citeproc{ref-beck-2021}{2021}) finds that during the pandemic, those with white collar occupations report to WfH for an average of 4.2 days compared to 1.5 for blue collar. This remains to be true in 2024 as Laß \& Wooden (\citeproc{ref-lass-2025}{2025}) find that WfH rates (defined here as working at least one full day at home per week) among employees are only pronounced in the white-collar job groups of managers (46\%), professionals (41.7\%), and clerical and administrative workers (36.3\%) while extremely uncommon in semi-skilled and unskilled blue-collar jobs (less than 1\%).

The nature of job also interacts with gender and education. As previously mentioned, WfH uptake is higher in females than males. Aside from conventional roles of women, the Australian Bureau of Statistics (2016a) demonstrates that 26\% of females while 15\% of males worked as Professionals; 22\% compared to 6\% of males worked as clerical and administrative workers. Akshay Vij (\citeproc{ref-vij-2023}{2023}) show the mentioned occupations to have the highest WfH uptake during the pandemic. Laß \& Wooden (\citeproc{ref-lass-2025}{2025}) explains that individuals who WfH tend to be well educated, as 65.8\% of such employees hold a university degree.

The DITRDCSA divides respondents (nonmanagerial and managerial employees who did not change their job during the pandemic) according to whether their location is a capital or a regional city. The eight regional cities is composed of Sydney, Melbourne, Adelaide, Perth, Brisbane, Canberra, Darwin and Hobart while Gold Coast-Tweed Heads, Newcastle-Maitland, Central Coast, Sunshine Coast, Wollongong, Geelong, Townsville, Cairns or Toowoomba make up the regional cities. They demonstrate that WfH capability and uptake are higher for the capital cities than the regional cities and this significant gap remains evident through early 2021. Among the non-managerial employees, WfH uptake and capability are higher in Central Business Districts (the five big capital cities) as well as in inner BCARR rings compared to their counterparts.

While firms often fear that restricting WfH may increase employee turnover, empirical evidence remains limited; using HILDA data, Mooi-Reci and Wooden (2025) find that greater WFH is associated with lower quit rates, though this effect is significant only among male employees (Laß \& Wooden (\citeproc{ref-lass-2025}{2025})). In US firms enforcing full return-to-office, many employees simply do not comply. Even with four-day office policies, compliance rates around 80\% still leave one or two members of a typical team working remotely on any given day, making fully in-person collaboration difficult. As a result, employees often report coming to the office only to spend most of the day on virtual calls. In response, many managers take no action---over 40\% do not enforce strict attendance, especially when work is still being completed. Some managers themselves prefer hybrid setups and quietly tolerate employees working from home more than required, reflecting widespread passive resistance to full return-to-office policies (Bloom (\citeproc{ref-bloom-2022}{2022})). In another study, ADP Research Institute (\citeproc{ref-adp-2022}{2022}) finds that 64\% of the workforce would consider looking for a new job if required to return to the office full time. Contrary to common assumptions, this reluctance is even higher among younger workers (aged 18--24). The study also reports that employees are increasingly willing to accept trade-offs for flexibility, with more than half indicating willingness to take a pay cut in exchange for a flexible work arrangements.

\section{Methodology}\label{methodology}

\subsection{Data source}\label{data-source}

\url{https://wfhresearch.com/wp-content/uploads/2022/10/Working-from-Home-Around-the-World-16-October-2022.pdf}

\section{Result}\label{result}

\section{Conclusion}\label{conclusion}

\clearpage
\vspace*{1cm}
\phantomsection

\section*{References}\label{references}
\addcontentsline{toc}{section}{References}

\vspace*{0.8cm}

\protect\phantomsection\label{refs}
\begin{CSLReferences}{1}{0}
\bibitem[\citeproctext]{ref-adp-2022}
ADP Research Institute. (2022). \emph{ADP research institute reveals pandemic-sparked shift in workers' priorities and expectations in new global study} {[}Press Release / Research Summary{]}. ADP, Inc. Retrieved from ADP, Inc. website: \url{https://mediacenter.adp.com/2022-04-25-ADP-Research-Institute-R-Reveals-Pandemic-Sparked-Shift-in-Workers-Priorities-and-Expectations-in-New-Global-Study}

\bibitem[\citeproctext]{ref-vij-2023}
Akshay Vij, H. B., Flavio F. Souza. (2023). Employee preferences for working from home in australia. \emph{Journal of Economic Behavior \& Organization}. \url{https://doi.org/10.1016/j.jebo.2023.08.020}

\bibitem[\citeproctext]{ref-barrero-2022}
Barrero, J. M., Bloom, N., \& Davis, S. J. (2022). \emph{Working from home around the world} (Working Paper No. 28745). National Bureau of Economic Research. Retrieved from National Bureau of Economic Research website: \url{https://www.nber.org/papers/w28745}

\bibitem[\citeproctext]{ref-beck-2021}
Beck, M. J., \& Hensher, D. A. (2021). Insights into working from home in australia in 2020: Positives, negatives and the potential for future benefits to transport and society. \emph{Transport Policy}, \emph{108}, 172--186. \url{https://doi.org/10.1016/j.tranpol.2021.05.007}

\bibitem[\citeproctext]{ref-bloom-2022}
Bloom, N. (2022). \emph{The great resistance: Getting employees back to the office} {[}Policy Brief{]}. Stanford Institute for Economic Policy Research (SIEPR). Retrieved from Stanford Institute for Economic Policy Research (SIEPR) website: \url{https://siepr.stanford.edu/publications/work/great-resistance-getting-employees-back-office}

\bibitem[\citeproctext]{ref-datasource-2022}
Cevat Giray Aksoy, N. B., Jose Maria Barrero, \& Zarate, P. (2022). \emph{Working from home around the world}. Retrieved from \url{https://wfhresearch.com/gswadata/}

\bibitem[\citeproctext]{ref-ceda-2024}
Committee for Economic Development of Australia (CEDA). (2024). \emph{Data insight: Working from home 2024} {[}Research Report{]}. CEDA. Retrieved from CEDA website: \url{https://cedakenticomedia.blob.core.windows.net/cedamediatest/kentico/media/ceda-landing-pages/241216_ceda-wfh-data-final.pdf}

\bibitem[\citeproctext]{ref-dit-2022}
Department of Infrastructure, Transport, Regional Development, Communications and the Arts (DITRDCSA). (2022). \emph{The role of socio-demographic and spatial characteristics in {Work from Home} in {Australia}} {[}Research Report{]}. Australian Government. Retrieved from Australian Government website: \url{https://www.infrastructure.gov.au/sites/default/files/documents/role-socio-demographic-spatial-characteristics-work-from-home-australia.pdf}

\bibitem[\citeproctext]{ref-lass-2025}
Laß, I., \& Wooden, M. (2025). Working from home: The australian experience. \emph{Australian Economic Review}, \emph{58}(2), 154--162. \url{https://doi.org/10.1111/1467-8462.70010}

\bibitem[\citeproctext]{ref-ruger-2024}
Rüger, H., Laß, I., Stawarz, N., \& Mergener, A. (2024). To what extent does working from home lead to savings in commuting time? A panel analysis using the australian HILDA survey. \emph{Travel Behaviour and Society}, \emph{37}, 100839. \url{https://doi.org/10.1016/j.tbs.2024.100839}

\end{CSLReferences}

\clearpage
\vspace*{1cm}

\phantomsection 
\appendix

\section*{Appendix}\label{appendix}
\addcontentsline{toc}{section}{Appendix}

\vspace*{0.8cm}

\setcounter{section}{0}
\renewcommand{\thesection}{A}

\setcounter{subsection}{0}
\renewcommand{\thesubsection}{A.\arabic{subsection}}

\numberwithin{table}{subsection}
\numberwithin{figure}{subsection}

\end{document}
